\section{The full integral cusp lattice}

% BEGIN M76L1
\paragraph{M76L1. The whole valuation lattice spans in positive weight.}
Fix the actual lower row satisfying the pins and the distinguished
uniformization laws used in ICL19. Let $K=c.F$, and let $C_w$ be its
original cusp space at a dominant weight $w$ with $w.n>0$. Define
\[
 L_w=\{f\in C_w:\ v(f)\ge0\text{ for every original valuation test}\}.
\]
Then $L_w$ is a finitely generated $\mathcal O_K$-module and
$K L_w=C_w$. This assertion concerns the whole set defined by the
valuation inequalities. The short sections used in its proof need
only generate the appropriate bounded sections over $K$.

Indeed ICL8 proves finite generation of $L_w$. For spanning, take
$f\in C_w$ and set $q=4[E:\mathbb Q]r+4$. By ICL10 the original
section $f^q$ is cuspidal of dominant theta exponent $qw.n\ge q$.
Boundedness carries to the principal charts by C76P2. Apply B76S2
and write $f^q=\sum_j a_jg_j$ with $a_j\in K$ and $g_j$ finitely
many identity-cover short sections. R76V4 gives $s_j\in K^\times$
such that $s_jg_j$ is nonnegative at every original test. There
is one nonzero $t\in\mathcal O_K$ for which $ta_j/s_j$ is integral
for every $j$: clear the denominators of this finite list.
Thus $tf^q$ is nonnegative at every test, by the valuation inequality
for sums. If $v$ lies over $\mathfrak p$, then
\[
 q\,v(tf)=(q-1)\operatorname{ord}_{\mathfrak p}(t)+v(tf^q)\ge0.
\]
This includes zero values, with valuation $+\infty$. Hence
$tf\in L_w$ and $f\in K L_w$. This applies unchanged to the zero
space and uses no common bound on the denominators as $g_j$ varies.
Under the inclusion of $C_w$ into its ambient rule space, the
definition of $\mathrm{latIn}$ gives precisely $L_w$ (ICL1).
Transport along that inclusion proves the corresponding
$\mathrm{FullLat}$ statement on the subtype.
% END M76L1

% BEGIN M76L2
\paragraph{M76L2. Integral lattices at both of the prescribed levels.}
Let $T:\mathrm{CoreJ}(D)$ satisfy its original $\mathrm{LawsJ}$
and let the given family of weights $\kappa$ satisfy $\mathrm{KapIs}$.
Then $\mathrm{LatFull}(T,\kappa)$ holds. Indeed for each original
row $j=i+1$ these hypotheses provide its pins and distinguished
uniformization, and $(\kappa_i).n=D.R.nV>0$ by ICL20.
Insert M76L1 on that row. For the canonical weight the same equality
of exponents and the defining row weight in ICL20 give the canonical
integral cusp-forms full-lattice assertion.

Now fix the original $\mathrm{FWTie}(T,R,\kappa,v_0)$ and any
admissible face $F$. Its injection from the cusp space identifies
the source lattice with the image in the prescribed rational range,
using its actual valuation and norm equalities (ICL21).
Consequently each stored row satisfies $\mathrm{MKFull}$.
For either Boolean level $b$ the identification of coordinates
carries this finite spanning lattice into the stored rational space.
The original lattice at that level is finite and spans, by the
transport in RTG22--RTG23. Thus $(R.\mathrm{toCx}(F)).\mathrm{MLat}(b)$
holds for both values of $b$, on the given realization. The valuation
argument proves the denominator needed here for each generator;
the stronger uniform estimate (ICL.14) on its individual torsion
sums remains a separate assertion.
% END M76L2
